Using LLMs in empirical SE research, whether as tools for researchers or for software engineers, offers several potential advantages but also raises fundamental challenges that cut across the study types described above.

\studytypeparagraph{Advantages}

The primary advantage of using LLMs for research tasks is \emph{speed, cost reduction, and scalability}. LLMs can annotate, judge, synthesize, and simulate faster and at lower cost than human researchers, with studies showing cost reductions of 50--96\% on various natural language tasks~\cite{DBLP:conf/emnlp/WangLXZZ21} (e.g., \citeauthor{DBLP:conf/chi/HeHDRH24} found that GPT-4 annotation required only two days and 122.08 USD compared to several weeks and 4,508 USD for a comparable MTurk pipeline~\cite{DBLP:conf/chi/HeHDRH24}). Similarly, augmenting or replacing human participants with LLM-generated virtual participants would reduce recruitment effort~\cite{DBLP:conf/vl/Madampe0HO24}. This efficiency unlocks scalability: qualitative research traditionally does not scale well to large samples, but LLMs let researchers process more text than human-only coding allows and support larger judgment datasets.
% GROUNDING DBLP:conf/emnlp/WangLXZZ21: "it costs 50% to 96% less to use labels from GPT-3 than using labels from humans"
% GROUNDING DBLP:conf/chi/HeHDRH24: "The main experiment with GPT-4 totaled only $122.08, compared to the $4,508 spent on MTurk"
% GROUNDING DBLP:conf/vl/Madampe0HO24: "our experience of recruiting software professionals as potential participants in our studies has met with fluctuating success ... This cost us more effort and time than anticipated."

LLMs can also \emph{automate} tasks such as coding qualitative data, assessing artifact quality, and generating summaries, reducing cognitive demands and resources required for qualitative research. Some studies suggest that LLMs may improve \emph{consistency}: ChatGPT's accuracy exceeded crowd workers by approximately 25\%, and LLMs can achieve higher inter-rater agreement than crowd workers and trained annotators~\cite{DBLP:journals/corr/abs-2303-15056}. However, both human and LLM judges exhibit systematic biases, such as preferring outputs with fake citations or rich formatting~\cite{DBLP:conf/emnlp/ChenCLJW24}, and no compelling evidence supports claims that LLMs are less biased than human annotators.
% GROUNDING DBLP:journals/corr/abs-2303-15056: "the zero-shot accuracy of ChatGPT exceeds that of crowd-workers by about 25 percentage points on average"
% GROUNDING DBLP:conf/emnlp/ChenCLJW24: "Results show that human and LLM judges are vulnerable to perturbations to various degrees, and that even the cutting-edge judges possess considerable biases."

LLMs could potentially provide \emph{access to otherwise-inaccessible research contexts}. \textit{If} virtual participants' behavior is sufficiently similar to human participants, LLMs could access underrepresented and hard-to-reach populations, strengthen generalizability and inclusiveness, impute missing data (see \synthesis), and enable research that is ethically problematic with real humans (e.g., questions that would force a human to relive past trauma do not harm an LLM).

Studying real-world LLM-based tools allows researchers to \emph{understand the state of practice}, uncovering usage patterns, adoption rates, and contextual factors, and \emph{generate hypotheses} about how LLMs affect developer productivity, collaboration, and decision-making. From an engineering perspective, developing LLM-based tools \emph{requires less task-specific engineering} than traditional SE approaches such as static analysis or symbolic execution, because a single model can handle diverse inputs without language-specific parsers or hand-crafted rules. Good benchmarks provide \emph{standardized evaluation and model comparison}, reducing research effort and supporting open science by providing common ground for sharing data and results. Benchmarks built for specific SE tasks can help identify LLM weaknesses and support optimization and fine-tuning.

\studytypeparagraph{Challenges}

The stochastic nature of LLM responses, where identical prompts may yield different outputs, \emph{undermines replicability} across all study types, complicating interpretation of experimental results and test-retest reliability. More broadly, \emph{reliability} is a persistent concern: Like any measurement instrument, \judges and \annotators should exhibit validity, test-retest reliability, inter-rater reliability, minimal error, and measurement invariance~\cite{DBLP:books/sp/24/RalphKAA24}. However, LLMs show substantial variability depending on the dataset and task~\cite{DBLP:journals/corr/abs-2306-00176, DBLP:journals/corr/abs-2406-18403}. They are sensitive to prompt variations~\cite{DBLP:journals/corr/abs-2304-11085} and option order~\cite{DBLP:conf/naacl/PezeshkpourH24}, can behave differently when reviewing their own outputs~\cite{DBLP:conf/nips/PanicksseryBF24}, and can be unreliable for high-stakes labeling~\cite{DBLP:conf/chi/Wang0RMM24}. Crucially, \emph{reliability does not imply validity}: a reliable LLM might be reliably inaccurate~\cite{DBLP:conf/coling/ZhouC024}. For tasks with no single correct answer, the statistical framework pushes outputs toward the most likely answer, which may not be the best one~\cite{DBLP:journals/corr/abs-2510-01171, DBLP:journals/corr/abs-2310-06452}. See \modelversion and \limitationsmitigations for reporting guidance on replicability and reliability.
% GROUNDING DBLP:books/sp/24/RalphKAA24: "Multi-item scales or multi-metric instruments and statistical measurement models are required to assess the degree to which instruments measure what they are supposed to measure (construct validity)."
% GROUNDING DBLP:journals/corr/abs-2306-00176: "LLM performance for text annotation is promising but highly contingent on both the dataset and the type of annotation task"
% GROUNDING DBLP:journals/corr/abs-2406-18403: "Models are reliable evaluators on some tasks, but overall display substantial variability depending on the property being evaluated."
% GROUNDING DBLP:journals/corr/abs-2304-11085: "even minor wording alterations in prompts or repeating the identical input can lead to varying outputs"
% GROUNDING DBLP:conf/naacl/PezeshkpourH24: "we demonstrate a considerable performance gap of approximately 13% to 75% in LLMs on different benchmarks, when answer options are reordered"
% GROUNDING DBLP:conf/nips/PanicksseryBF24: "One such bias is self-preference, where an LLM evaluator scores its own outputs higher than others' while human annotators consider them of equal quality."
% GROUNDING DBLP:conf/chi/Wang0RMM24: "they are often erroneous for complex or domain-specific tasks and may introduce bias when compared to human annotations"
% GROUNDING DBLP:conf/coling/ZhouC024: "GPT-3.5 can understand human-written reviews, distinguish emotions, and give consistent scores"
% GROUNDING DBLP:journals/corr/abs-2510-01171: "Post-training alignment often reduces LLM diversity, leading to a phenomenon known as mode collapse."
% GROUNDING DBLP:journals/corr/abs-2310-06452: "RLHF significantly reduces output diversity compared to SFT across a variety of measures."

LLMs and LLM-based tools \emph{evolve rapidly}, and so do practitioners' workflows around them. This combination complicates longitudinal comparisons and can quickly render study findings obsolete.
Findings tied to a specific model version may not transfer to later releases even within the same model family, complicating reproducibility and longitudinal comparisons.

The prevalence of \emph{proprietary tools and opaque training data} limits researchers' ability to assess and mitigate biases, and enables benchmark contamination~\cite{DBLP:journals/corr/abs-2410-16186}. LLMs exhibit \emph{bias} in multiple forms: tendencies to overestimate certain labels~\cite{DBLP:journals/corr/abs-2304-10145}, well-documented fairness issues~\cite{DBLP:journals/coling/GallegosRBTKDYZA24}, and the potential to reinforce prejudices. When simulating human participants, LLMs are \enq{likely to both misportray and flatten the representations of demographic groups}~\cite{DBLP:journals/natmi/WangMD25}. See \openllm and \limitationsmitigations for mitigation strategies.
% GROUNDING DBLP:journals/corr/abs-2410-16186: "LLM creators do not always disclose the exact details of the datasets used, and it is plausible that the model is trained on benchmark datasets intentionally or unintentionally"
% GROUNDING DBLP:journals/corr/abs-2304-10145: "for stance detection, ChatGPT may have a tendency to overestimate certain stances, emphasizing the need for human moderation of labels"
% GROUNDING DBLP:journals/coling/GallegosRBTKDYZA24: "Despite this success, these models can learn, perpetuate, and amplify harmful social biases"
% GROUNDING DBLP:journals/natmi/WangMD25: "We argue analytically for why LLMs are likely to both misportray and flatten the representations of demographic groups"

\emph{Evaluation difficulty} is inherent in studying LLM-based tools and benchmarks. Benchmarks may lack construct validity~\cite{DBLP:books/sp/24/RalphKAA24}, usually do not capture the full complexity of software engineering work~\cite{Chandra2025benchmarks}, and may lead to \emph{overconfidence} and overfitting~\cite{DBLP:journals/corr/abs-2412-03597}. LLMs are also susceptible to adversarial manipulation where semantics-preserving changes may deceive them into accepting flawed artifacts~\cite{DBLP:journals/corr/abs-2510-24367}. See \benchmarksmetrics for detailed guidance on benchmark design, including~\citeauthor{cao2025should}'s guidelines for coding task benchmarks~\cite{cao2025should}.
% GROUNDING DBLP:books/sp/24/RalphKAA24: "We rarely see statistical measurement models in repository mining, benchmarking, or other lab-based quantitative research."
% GROUNDING Chandra2025benchmarks: "Does the suite represent the software engineering work, including its complexity and diversity, that we want AI to help with in the first place?"
% GROUNDING DBLP:journals/corr/abs-2412-03597: "models excel at standardized tests while failing to demonstrate genuine language understanding and adaptability"
% GROUNDING DBLP:journals/corr/abs-2510-24367: "These attacks subtly alter software artifacts to deceive the judge into accepting flawed or malicious code."
% GROUNDING cao2025should: "we introduce a comprehensive code benchmark guideline, HOW2BENCH, comprising 55 rigor-oriented checklists, each annotated with its relative threat level to the benchmark validity"

A fundamental challenge is \emph{philosophical and methodological incongruence} with qualitative research. In a recent open letter, 419 experienced qualitative researchers argued \enq{that analytic approaches such as reflexive thematic analysis are human research practices requiring a subjective, positioned, and reflexive researcher and therefore the use of GenAI in such approaches is not methodologically congruent}~\cite{jowsey2025reject}. LLMs lack the capacity for genuinely reflexive qualitative analysis because they operate on statistical prediction without understanding the meaning of the data being analyzed~\cite{DBLP:journals/corr/abs-2306-13298, DBLP:journals/corr/abs-2510-18456}. Effective use requires structured prompts and careful human oversight~\cite{DBLP:conf/wsese/BarrosANKKNB25}. See \humanvalidation and \limitationsmitigations for guidance.
% GROUNDING jowsey2025reject: "analytic approaches such as reflexive thematic analysis are human research practices requiring a subjective, positioned, and reflexive researcher and therefore the use of GenAI in such approaches is not methodologically congruent"
% GROUNDING DBLP:journals/corr/abs-2306-13298: "Despite LLMs' adeptness at generating linguistically coherent responses, they do not genuinely comprehend the meanings, nuances, and deeper implications of words and phrases."
% GROUNDING DBLP:journals/corr/abs-2510-18456: "LLMs can surface patterns, candidate framings, and alternative readings, but they do not make meaning in the epistemological sense; that remains the role of reflexive human interpretation."
% GROUNDING DBLP:conf/wsese/BarrosANKKNB25: "they are highly dependent on well-structured prompts, which, if poorly designed, can lead to inaccurate or incomplete outputs"

There is \emph{insufficient evidence} for the effectiveness of LLMs in most research roles. No compelling evidence exists that LLMs can accurately simulate human participants~\cite{DBLP:journals/corr/abs-2508-06950} or reliably judge most relevant properties of SE artifacts. Gathering such evidence for each specific usage may be quite difficult~\cite{DBLP:journals/ais/HardingDLL24}. LLMs may be better suited for augmenting rather than replacing human researchers~\cite{DBLP:conf/emnlp/WangLXZZ21, DBLP:conf/chi/HeHDRH24}, but even then, limited evidence exists that augmentation increases \emph{effectiveness} as well as \emph{efficiency}. When LLM outputs are incorrect, they can negatively \emph{affect human judgment}~\cite{DBLP:conf/www/HuangKA23a}. See \humanvalidation for mitigation strategies.
% GROUNDING DBLP:journals/corr/abs-2508-06950: "We conclude that LLMs do not simulate human psychology and recommend that psychological researchers should treat LLMs as useful but fundamentally unreliable tools"
% GROUNDING DBLP:journals/ais/HardingDLL24: "In scenarios like this, the informativeness of the LLM's output is impossible to assess without doing further confirmatory work with human participants."
% GROUNDING DBLP:conf/emnlp/WangLXZZ21: "we propose a novel framework of combining pseudo labels from GPT-3 with human labels, which leads to even better performance"
% GROUNDING DBLP:conf/chi/HeHDRH24: "the research focus might shift from 'using AI to support human labelers' to 'using humans to enhance AI labeling.'"
% GROUNDING DBLP:conf/www/HuangKA23a: "In a case when ChatGPT's decision is wrong, such capability might be a danger to misleading lay people and potentially constructing wrong labels."

Majority voting across multiple outputs improves reliability~\cite{DBLP:journals/corr/abs-2304-11085} but increases \emph{cost and environmental impact}. While open models are available, the most capable ones require substantial hardware; relying on \emph{cloud-based APIs} introduces concerns related to data privacy, security, and replicability. See \limitationsmitigations for sustainability considerations.
% GROUNDING DBLP:journals/corr/abs-2304-11085: "majority decisions on repeated classification outputs of the same input led to improvements in consistency"

Field study findings face \emph{generalizability} challenges because outcomes may be highly sensitive to individual differences, usage patterns, goals, and contexts. Field studies must be \enq{dependable}~\cite{Sullivan2011-ub} beyond traditional validity criteria, which complicates methodology; see \limitationsmitigations for a detailed threat taxonomy.
% GROUNDING Sullivan2011-ub: "Dependability is gained though consistency of data, which is evaluated through transparent research steps and research findings."
